% The template needs to be compiled using LuaLaTeX or XeLaTeX

\documentclass{article}
\usepackage{jtdh}
% \definecolor{Mycolor2}{HTML}{00F9DE}
\usepackage[utf8]{inputenc} % allow utf-8 input
\usepackage[slovene]{babel} % Change to english if english paper
\usepackage{url}   % simple URL typesetting
\AtBeginDocument{\urlstyle{APACsame}}
\usepackage[hang,flushmargin]{footmisc} % no indentation in footnotes
\usepackage{hyperref}
\hypersetup{colorlinks = true, linkcolor = black, urlcolor = gray, citecolor = black, anchorcolor = black}
\usepackage{booktabs}       % professional-quality tables
\usepackage{amsfonts}       % blackboard math symbols
\usepackage{nicefrac}       % compact symbols for 1/2, etc.
\usepackage{xcolor}
\usepackage{graphicx}
\usepackage{tabularx}
\usepackage[backend=biber, style=apa, uniquelist=false]{biblatex}
\addbibresource{references.bib}

\usepackage[]{gb4e}  %for glossed examples
\let\eachwordone=\sffamily %or \normalfont
\let\eachwordtwo=\sffamily  

\usepackage{enumitem}
\graphicspath{{img/}}
% table styles
\aboverulesep=0ex
\belowrulesep=0ex

%\addbibresource{references.bib} %Imports bibliography file 

\title{Naslov prispevka}

\author{
  Ime PRIIMEK,\textsuperscript{1}\orcidlink{0000-0000-0000-0000} Ime PRIIMEK\textsuperscript{1, 2}\orcidlink{0000-0000-0000-0001}
}
\institutions{
  \textsuperscript{1}Ustanova avtorja. Primer ustanove je: Inštitut za novejšo zgodovino, Ljubljana, Slovenija \rorlink{04wr6tn02}\par
  \textsuperscript{2}Ustanova avtorja
}

\begin{document}
\maketitle

\begin{abstract}
Priporočeni obseg izvlečka je od 10 do 15 vrstic. Izvlečku sledi od 3 do 5 ključnih besed v slovenščini (oz. jeziku, v katerem je prispevek napisan).
\end{abstract}
\keywords{ključna beseda 1, ključna beseda 2, ključna beseda 3}

\section{ORCID in ROR}

Ob imenu vsakega avtorja, za nadpisanim označevalcem ustanove, je naveden trajni identifikator \href{https://orcid.org/}{ORCID} (Open Researcher and Contributor ID). Ne pozabite ORCID-ikone povezati s pravilnim URL-naslovom. Če ORCID-a nimate, se lahko registrirate ali pa s seznama avtorjev ikono ORCID odstranite.

Na koncu vsake vrstice z navedbo institucije je naveden trajni identifikator raziskovalnih ustanov \href{https://ror.org/}{ROR} (Research Organization Registry). Ne pozabite ROR-ikone povezati s pravilnim URL-naslovom. Če institucija, ki ste ji pridruženi, ni vpisana v ROR, pri tej instituciji ikono ROR odstranite.

\section{Številčenje naslovov, črkopis, velikost črk, razmik, poravnava, obseg}

Številčenje naslovov in podnaslovov je desetiško in se vedno začne z 1 (1, 1.1, 1.1.1; zaželeno je, da členitev nima več kot treh nivojev). Priporočena dolžina prispevka je navedena na spletni strani konference.

\section{Navajanje virov, sprotne opombe, tabele, slike}

\subsection{Navedba avtorja}

Navajanje virov v besedilu naj sledi 7.\ izdaji \href{https://apastyle.apa.org/}{slogovnega priročnika APA}. 
Uporabite naslednji obliki navedkov v besedilu:
\begin{itemize}
    \item celoten navedek v oklepaju (ang. \textit{parenthetical citation}): \parencite{leech1992corpora};
    \item priimek izven oklepaja, letnica izdaje v oklepaju (ang. \textit{narrative citation}): \textcite{biber1998corpus} navajajo\ldots{}
\end{itemize}
 
V primeru dobesednega navedka se doda številko strani: \parencite[107]{leech1992corpora} ali \parencite[107--108]{leech1992corpora}. 

Dela z dvema avtorjema se navaja z obema avtorjema: \parencite{studije2005korpusnem}. Po standardu APA \parencite[gl.][]{apa-author-date} se znak \& uporablja v primerih, kjer je celoten navedek v oklepaju; v navedku, kjer je  priimek izven oklepaja, letnica izdaje pa v oklepaju, pa se jih zapiše tako: \textcite{studije2005korpusnem} pokažejo, da \ldots{}    

Dela s tremi ali več avtorji se navaja le s prvim avtorjem: \parencite{biber1998corpus}, pri čemer se v razdelku Literatura navede vse avtorje. 

Več del, navedenih v enem navedku v oklepaju, naj bo razvrščenih po abecednem vrstnem redu in ločenih s podpičji: \parencite{erjavec-infotheca, tei}. 

Dela istega avtorja (ali istih avtorjev) iz istega leta se razlikuje z dodajanjem črk (a, b, c itd.) k letu objave: \parencite{erjavec-multext}. 

Če bi dve ali več del iz istega leta vodili do enakih skrajšanih navedkov (npr. \textit{Krek in sod., 2020}), jih razlikujte z dodajanjem črk k letu, na primer:  \parencite{krek2020a-ssj500k, krek2020b-gf} namesto navajanja dodatnih imen avtorjev. (To je izjema od sloga APA, 7.\ izd., uvedena zaradi poenostavitve navajanja virov v zborniku.)

Nekaj dodatnih zgledov citiranja:

    \begin{itemize}
        \item Knjiga: \parencite{biber1998corpus, studije2005korpusnem};
        \item Članek v znanstveni reviji: \parencite{biber1996investigating, erjavec-infotheca};
        \item Članek v monografiji ali konferenčnem zborniku: \parencite{leech1992corpora, erjavec-multext};
        \item Slovar: \parencite{sinclair1987collins, sskj};
        \item Spletna stran (ang. \textit{webpage}): \parencite{scott2008wordsmith, tei}.
    \end{itemize}

Kadar se sklicujete na spletno mesto (ang. \textit{website}) kot celoto in ne na posamezno spletno stran, v besedilu navedite ime spletnega mesta in mu neposredno dodajte hiperpovezavo, npr. \href{http://creativecommons.org/}{Creative Commons} \parencite{apa-website}.

Če manjka datum, navedite vir tako: \parencite{slobench-ner}.

Če je delo dostopno na spletu, je v Literaturi obvezno navesti njegov URL, pri čemer naj bo URL dodan kot hiperpovezava. Če obstaja, naj se kot povezavo navede trajni identifikator (DOI ali Handle URL). Hiperpovezave, ki niso trajni identifikatorji, imajo naveden datum dostopa v obliki: Pridobljeno 23.\ januarja 2012, \url{http://bos.zrc-sazu.si/sskj.html}. Vsako enoto v seznamu literature zaključuje pika, razen če se konča s hiperpovezavo.

Za zapis številčnega razpona v Literaturi vedno uporabimo pomišljaj in ne vezaj. Na primer, >>str.\ 110--112<< je pravilen zapis, medtem ko zapis >>str.\ 110-112<< s pravopisnega vidika ni pravilen. 

\subsection{Citiranje jezikovnih virov}

Jezikovni viri, kot so korpusi, računalniški leksikoni, jezikovni modeli itd., se navajajo v razdelku za literaturo, in sicer na enak način kot knjige, kar pomeni, da so navedeni tudi avtorji. Repozitorij, kjer je korpus deponiran, z vidika citiranja ustreza založniku pri knjigah. Na primer: \textcite{parliamentaryCorpus}. 

Če se isti vir nahaja na več lokacijah, navedemo različico, ki smo jo dejansko uporabili. Na primer: \parencite{gigafida-clarin} v primeru korpusa \textit{Gigafida 2.0}, če ste ga uporabili preko konkordančnika CLARIN.SI, oziroma \parencite{gigafida-cjvt}, če ste ga uporabili preko konkordančnika CJVT.

Če za vir obstaja trajni identifikator, ga navedemo: \url{http://hdl.handle.net/11356/1748} pri korpusu \textit{siParl 3.0} \parencite{parliamentaryCorpus}.\footnote{Velja opozoriti, da \url{https://www.clarin.si/repository/xmlui/handle/11356/1748} ni trajni identifikator.}

\subsection{Sprotne opombe, izpusti in daljši navedki}

Številko za opombo vnesemo za ločilom,\footnote{To je primer za besedilo opombe.} tudi v primeru seznama avtorjev na začetku dokumenta. Če je v opombi samo URL-naslov, na koncu ni pike.

Izpust iz navedka označimo z oglatimi oklepaji in tremi pikami: [\ldots]. Na začetku in na koncu navedka oznaka izpusta ni potrebna. Daljši navedki (več kot 3 vrstice) naj bodo postavljeni v samostojen odstavek z odmikom od levega roba. Taki navedki so brez narekovajev. Pred umaknjenim odstavkom ni prazne vrstice, za njim prav tako ne. Primer daljšega navedka:

\begin{quote}
Da se je Ramovš zavzel za vpeljavo teh oblik –- četudi pogojno –- v oficielni pravopis, pa je bila posledica njegovega novega, samostojnega študija slovenskega jezika, pri katerem je odkril takoimenovani kratki nedoločnik kot prastaro, še predtrubarsko varianto dolgega nedoločnika \parencite[198]{vodusek1959}.
\end{quote}

Narekovaji naj bodo dvojni srednji. Na primer: \citeauthor{polajnar2021samoreference} trdita, da je za >>akademski diskurz [\ldots] tradicionalno poleg intersubjektivnosti značilna poudarjena intelektualizacijska vloga<< \parencite*[413]{polajnar2021samoreference}. Če se poved ne konča s citatom, naj bo končno ločilo postavljeno pred končno navednico.

\subsection{Navedba tabel in slik}

\subsubsection{Številčenje in naslavljanje tabel ter slik}

Tabele in slike številčimo zaporedno. Nanje se sklicujemo z rabo velike začetnice (npr. v Tabeli \ref{slogi} prikazujemo …). Naslov je nad tabelo ali sliko, zaključuje ga pika. Besedilne vrednosti so v tabeli poravnane levo, številčne pa desno. Za morebitno ločevanje tisočic rabimo piko, za decimalke pa vejico.

\begin{table}[h]
\caption{Naslov tabele.}
\label{slogi}
\begin{tabularx}{\textwidth}{X|X|r}
\toprule
\textit{Kategorija 1} & \textit{Kategorija 2} & \textit{Kategorija 3 –   številčne vrednosti} \\ \midrule
Prva vrstica          & Besedilo tabele       & 1.900,12                                      \\
Druga vrstica         & Besedilo tabele       & 13.934,34                                     \\ \bottomrule
\end{tabularx}
\end{table}

\begin{figure}[h]
    \centering
    \caption{Opis slike se nahaja nad sliko in se konča s piko.}
    \includegraphics[width=0.7\linewidth]{img/image-dog.jpg}
    \label{fig:placeholder}

\figurenote{\textit{Opomba}. Vir: Lava [fotografija], Denali National Park and Preserve, 2013, Flickr (\href{https://www.flickr.com/photos/denalinps/8639280606/}{https://www.flickr.com/photos/denalinps/8639280606/}). CC BY 2.0.\nocite{denali2013}}

\end{figure}

\subsubsection{Navajanje slik, ki zahtevajo priznanje avtorstva}

Fotografije, ki zahtevajo priznanje avtorstva, kot v Sliki \ref{fig:placeholder}, morajo biti pravilno navedene. V skladu s smernicami APA \parencite{apa-clipart} je treba v opombi pod sliko navesti: naslov dela, opis v oglatih oklepajih, avtorja, leto objave, vir (ime spletnega mesta in URL) in ime licence \parencite{apa-clipart} -- to se uporabi namesto navedka v besedilu.

Referenco za sliko je treba vključiti tudi v seznam literature na koncu. Vnos v seznamu literature vsebuje enake elemente, kot so v opombi k sliki (razen licence), toda v drugačnem vrstnem redu: avtor, leto objave, naslov, opis v oglatih oklepajih in vir (ime spletnega mesta in URL) \parencite{apa-clipart}. Uporabniki LaTeX-a naj za prikaz takšnega vnosa v seznamu literature uporabijo ukaz \verb|\nocite{...}|.

Na prosojnicah oziroma drsnicah sta številka in naslov slike neobvezna, opomba s pripisom avtorstva pa je obvezna \parencite{apa-clipart}.

\subsection{Posebni slogi}

Priporočeni slog za alineje je naslednji:
\begin{itemize}
    \item seznam, postavka 1;
    \item seznam, postavka 2.
\end{itemize}




\acknowledgment{
Primer besedila v zahvali.
}



\urlstyle{same}
\printbibliography


\newpage
\titleeng{Naslov članka v angleščini}
\begin{abstracteng}
Na lastni strani na koncu prispevka je naslov članka in povzetek v angleškem jeziku (oz. naslov in povzetek v slovenskem jeziku, če je članek napisan v angleščini). Povzetku sledi tri do pet ključnih besed. Na isti strani (na dnu) je obvestilo o licenci.
\end{abstracteng}

\keywordseng{keyword 1, keyword 2, keyword 3}

\end{document}